% ---------------------------------------------------------------------------
% Response to R2 Major 6 -- JEDC revision
% Body only: paste after the previous comment in the combined response file.
% Needs booktabs in the preamble. The top heading is \section (numbered) and
% the subheadings are \subsection*, matching the author's edit.
% ---------------------------------------------------------------------------

\section{R2, Major 6}
{\itshape The magnitude of the model's spread response}

\begin{quote}\itshape
The effect in the model is much smaller than that in the data [\ldots] A
discussion of what determines this magnitude in the model would help the reader
interpret the result. Also, the author could estimate regression (1) on
model-simulated data, with volatility in the same units as the empirical
measure, so that the model's coefficients can be compared with the empirical
ones.
\end{quote}

Regression (1) has been estimated on simulated data as suggested. The following
table shows the descriptive statistics of the empirical volatility vs the
simulated volatility from the model.

\begin{table}[h]
\centering\footnotesize
\begin{tabular}{lccc}
\toprule
 & min & mean & max \\
\midrule
data, the paper's volatility series & 8.86 & 12.90 & 31.58 \\
model, $100\exp(\text{state}/2)$    & 8.77 & 13.79 & 20.64 \\
\bottomrule
\end{tabular}
\end{table}

The one asymmetry is at the top: the data reach 31.6 and the model's grid stops
at 20.6, so the model never experiences a quarter as volatile as 2008Q4. That
matters below.

Regression (1) in the paper relates the contemporaneous change in the log spread
to the \emph{previous} quarter's change in the log oil price, which avoids the
simultaneity between the two within a quarter. The simulated counterpart does
the same thing. It also controls for the change in debt to GDP:
\texttt{dy\_diff} is the first difference of \texttt{dy3}, which is commercial
public and publicly guaranteed external debt as a percent of GDP. The simulated
regression matches both choices, with the simulated debt stock taken over four
times quarterly output so the ratio is on an annual-GDP basis, averaging 26.1\%
against 0 to 25.3\% for the data series.

\begin{table}[h]
\centering\footnotesize
\begin{tabular}{lcc}
\toprule
 & model & data \\
\midrule
oil price volatility               & $+0.0047$ & $+0.0100$ \ ($t$ 2.29) \\
oil price change, previous quarter & $+0.0007$ \ ($t$ 0.4) & $+0.1416$ \ ($t$ 0.93) \\
change in debt / GDP               & $+0.1338$ \ ($t$ 297) & $+0.0027$ \ ($t$ 0.49) \\
\bottomrule
\end{tabular}
\end{table}

Standard errors are two-way clustered by country and quarter in the data,
following Major 1c of the other report, and clustered by simulated history in
the model, where each of the 500 independent histories plays the role a country
plays in the panel.

\textbf{On volatility, which is the paper's subject, the model and data both
show a positive and statistically significant coefficient and the model delivers
47\% of the empirical coefficient.}

\textbf{On the oil price change the model and the data agree, and they agree on
zero.} The simulated coefficient is $+0.0007$ with a $t$ of 0.4 and the
empirical one $+0.1416$ with a two-way $t$ of 0.93. The 6 to 8 basis points in
the comment are computed from the empirical coefficient of 0.14, whose standard
error covers zero.

The model's spread responds to the oil price in the same quarter rather than the
next: on the contemporaneous change the coefficient is $+1.03$. It then retraces
as the price mean-reverts, which is why the previous quarter's change shows
nothing.

\textbf{On debt the model is far more responsive than the panel}, $+0.134$
against $+0.0027$, and the difference looks like measurement rather than
mechanism. In the model debt is the state variable and the spread is a function
of it, so the relationship is exact. In the data the regressor is one debt
aggregate, reported annually and carried across the quarters of each year, which
is a noisy and slow-moving proxy for the leverage a lender prices; an attenuated
coefficient is what that predicts, and the empirical estimate is indeed
indistinguishable from zero.

\subsection*{What governs the magnitude in the model}

Three things, in order of how much they bind.

\textbf{The size of the shock.} The oil bill is 9 to 10\% of GDP and the
substitution elasticity is below one, so the quantity response does not offset
the price and a 10\% oil price rise is close to a 1\% of GDP resource loss. The
model turns that into an 11\% rise in the spread on impact. The transmission is
not weak; the shock is small relative to what usually moves sovereign spreads.

\textbf{Where the volatility grid stops.} The model's volatility reaches 20.6 in
the empirical units against 31.6 in the data, so the states that would generate
the largest responses are outside its support. This is a direct and quantifiable
reason its volatility coefficient is the smaller of the two.

\textbf{Risk-neutral lenders.} Spreads in the model are pure expected-default
compensation, with no term for the price of risk. Major 5 of the other report
asks about precisely this, and our reply there reports what an exogenous
stochastic discount factor does. That is the most promising route to closing the
remaining gap on the volatility channel.

\subsection*{Suggested manuscript wording}

\begin{quote}
Regression (1) estimated on model-simulated data, with volatility expressed in
the units of the empirical measure and the same controls as in the panel, gives
a volatility coefficient of 0.0047 against 0.0100, so the model reproduces about
47 per cent of the estimated response. On the lagged oil price change the
simulated coefficient is 0.0007 and the empirical one 0.14, and neither is
distinguishable from zero once standard errors are clustered by quarter as well
as by country; the model's spread responds to the oil price within the quarter
rather than after it, with a contemporaneous coefficient of 1.03. The magnitude
of the model's response is governed by the size of the shock -- a 10 per cent
oil price change is close to a 1 per cent of GDP resource loss at an oil bill of
9 to 10 per cent of GDP -- by the range of the volatility grid, which stops
short of the most volatile quarters in the data, and by the assumption of
risk-neutral lenders, which leaves spreads as pure expected-default
compensation.
\end{quote}
